% !TEX program = xelatex
% Jiarui Zhang - Resume for applying to quantitative research jobs (Chinese)
\documentclass[11pt,a4paper]{article}

\usepackage[left=0.68in,right=0.68in,top=0.55in,bottom=0.55in]{geometry}
\usepackage{fontspec}
\usepackage{xeCJK}
\usepackage{xcolor}
\usepackage{titlesec}
\usepackage{enumitem}
\usepackage{tabularx}
\usepackage[hidelinks]{hyperref}
\usepackage{microtype}

\setmainfont{TeX Gyre Termes}
\setsansfont{TeX Gyre Heros}
\setCJKmainfont{Noto Serif CJK SC}
\setCJKsansfont{Noto Sans CJK SC}
\newfontfamily\latinGaramond{EB Garamond}
\setCJKfamilyfont{kai}{AR PL UKai CN}[BoldFont={AR PL UKai CN}]
\definecolor{accent}{HTML}{0A6670}
\definecolor{textgray}{HTML}{222222}

\pagestyle{empty}
\setlength{\parindent}{0pt}
\setlength{\emergencystretch}{2em}
\titleformat{\section}{\large\bfseries\sffamily\color{accent}}{}{0pt}{}[\vspace{-3pt}\color{accent}\rule{\textwidth}{0.55pt}]
\titlespacing*{\section}{0pt}{5pt}{3pt}
\setlist[itemize]{leftmargin=1.25em,itemsep=0.5pt,topsep=1pt,parsep=0pt,partopsep=0pt}

\newcommand{\entry}[4]{%
  \textbf{#1} \hfill {\small #2}\\[-1pt]
  #3 \hfill {\small #4}%
}

\newcommand{\weblink}[2]{%
  \href{#1}{\textcolor{accent}{\underline{#2}}\,{\footnotesize$\nearrow$}}%
}

\begin{document}

\begin{center}
  {\LARGE\bfseries\color{textgray} {\CJKfamily{kai}张家瑞}}\\[3pt]
  \href{mailto:jiarui-z23@mails.tsinghua.edu.cn}{jiarui-z23@mails.tsinghua.edu.cn}
  \quad $\vert$ \quad +86 189 2512 1908
  \quad $\vert$ \quad \href{https://ezoicoder.github.io}{ezoicoder.github.io}
\end{center}

\vspace{-3pt}
\textbf{个人概述：}我是清华大学姚班四年级本科生，主要关注理论计算机科学与机器学习，研究经历涵盖随机采样、电路复杂度与高效机器学习系统。

\section*{教育背景}
\textbf{清华大学}，北京 \hfill {\small 2023 年 9 月—2027 年 6 月（预计）}\\[-1pt]
计算机科学与技术工学学士，交叉信息研究院\\[-1pt]
GPA: 3.8/4.0 \quad $\vert$ \quad TOEFL iBT: 5.5/6.0\\[-1pt]
\textbf{核心课程：}算法设计（A-）、计算理论（A）、抽象代数（A）、自然语言处理（A+）

\section*{荣誉}
\begin{itemize}
  \item \textbf{2021 年全国青少年信息学奥林匹克竞赛（NOI 2021）}——金牌
  \item \textbf{NOI 2022 广东省队选拔}——第 1 名
\end{itemize}

\section*{经历}
\entry{本科生研究}{2026 年 8 月—至今}{清华大学张景昭教授指导}{}
\begin{itemize}
  \item 为正则语言和马尔可夫链构造 randomized-$AC^0$ 采样器。
  \item 证明 revision 能够降低扩散语言模型的串行计算量。
\end{itemize}

\entry{应用研究实习生}{2025 年 7 月—2025 年 9 月}{腾讯 CDG}{金融大语言模型}
\begin{itemize}
  \item 为金融大语言模型设计模型打分与提示词驱动的奖励评分方法，围绕引用有效性、内容准确性、时效性与指令遵循能力构建评测流程。
  \item 将引用格式准确率由 40\% 提升至 98\%，内容与时效性准确率由 78\% 提升至 97\%，IFEval 准确率由 38\% 提升至 80\%。
\end{itemize}

\entry{研究实习生}{2026 年 1 月—2026 年 8 月}{袁彬航课题组，香港科技大学}{高效机器学习系统}
\begin{itemize}
  \item 开发面向 Tree Attention 的分布式训练方法，降低显存占用并提升多轮大语言模型训练吞吐量。
  \item 参与 $D^2$SD 双扩散草稿模型研究，通过推测解码加速大语言模型推理。
\end{itemize}

\section*{论文发表}
\begin{itemize}
  \item \textbf{Jiarui Zhang*}, Yuchen Yang*, Ran Yan*, Zhiyu Mei, Liyuan Zhang, Daifeng Li, Wei Fu, Jiaxuan Gao, Shusheng Xu, Yi Wu, Binhang Yuan. \href{https://arxiv.org/abs/2602.00482}{``AREAL-DTA: Dynamic Tree Attention for Efficient Reinforcement Learning of Large Language Models.''} \textbf{ICML 2026}.
  \item Liyuan Zhang, \textbf{Jiarui Zhang}, Jinwei Yao, Ran Yan, Yuchen Yang, Jiahao Zhang, Tongkai Yang, Yi Wu, Binhang Yuan. \href{https://arxiv.org/abs/2606.04446}{``$D^2$SD: Accelerating Speculative Decoding with Dual Diffusion Draft Models.''} arXiv:2606.04446.
  \item \textbf{Jiarui Zhang*}, Chengwei Liang*, Haozhe Jiang, Binhang Yuan, Jingzhao Zhang. \href{https://zenodo.org/records/23105365}{``Revision Provably Reduces Sequential Computation in Diffusion Language Models.''} Zenodo preprint, 2026.
\end{itemize}
{\footnotesize * 共同一作。}

\section*{兴趣爱好}
\begin{itemize}
  \item 毽球，清华大学毽球队队员
  \item 徒步（喜欢单日爬升约 $1000\,\mathrm{m}$ 的路线）
\end{itemize}

\end{document}
